\clearpage
\thispagestyle{empty}
\pdfbookmark[0]{Magnetische Beobachtungen. März 1835. Ausgabe und Quellen}{march1835-edition}
\begin{center}\Large\textbf{Magnetische Beobachtungen\\März 1835}\end{center}

Der vollständige Bericht vom März 1835 über Carl Friedrich Gauß’ magnetische Beobachtungen liegt hier als KI-Übertragung in modernes Deutsch mit KI-erstelltem Satz vor. Grundlage ist die Fassung in den \textit{Werken}, Band V (Göttingen, 1867), S.~528--536, einschließlich der sechs Tabellenzeilen mit den Monatsmitteln und der Fußnote zur Aufhängung des Stabes. Der Bericht gibt einen Auszug aus Gauß’ Vortrag vor der Königlichen Sozietät in Göttingen wieder und spricht von ihm in der dritten Person.

Der Erstdruck erschien in den \textit{Göttingischen gelehrten Anzeigen}, Nr.~36, 7. März 1835, S.~345--357. Johann Christian Poggendorffs Zeitschrift druckte ihn als Beitrag XVI in den \textit{Annalen der Physik und Chemie}, Band~34 (110), 1835, S.~546--556, mit eigenen redaktionellen Anmerkungen und Verweisen nach.

\medskip\noindent
\href{https://gdz.sub.uni-goettingen.de/id/PPN236006339}{Deutsche Werkausgabe (1867)}\par
\noindent\href{https://www.digitale-sammlungen.de/view/bsb10538708?page=351}{Erstdruck (1835)}\par
\noindent\href{https://www.digitale-sammlungen.de/view/bsb10130347?page=558}{Nachdruck in den Annalen (1835)}
\clearpage
